% !PRINTED 77
\weberpage{77}
%<<CONT
\weq{2}{\vartheta_{01}^2\,\vartheta_{01}(u+v)\,\vartheta_{01}(u-v) = \vartheta_{01}^2(u)\,\vartheta_{01}^2(v) - \vartheta_{11}^2(u)\,\vartheta_{11}^2(v).}
\weq{3}{\vartheta_{10}^2\,\vartheta_{10}(u+v)\,\vartheta_{10}(u-v) = \vartheta_{10}^2(u)\,\vartheta_{10}^2(v) - \vartheta_{11}^2(u)\,\vartheta_{11}^2(v).}
\weq{4}{\vartheta_{01}^2\,\vartheta_{11}(u+v)\,\vartheta_{11}(u-v) = \vartheta_{11}^2(u)\,\vartheta_{01}^2(v) - \vartheta_{01}^2(u)\,\vartheta_{11}^2(v).}

Diese vier Gleichungen können, wie schon erwähnt, durch Anwendung der Formeln (13) des vorigen Paragraphen in mannigfacher Weise umgeformt werden. Die folgenden sechs Formeln haben nur \gesperrt{eine} Form. Es ist, um wieder mit einem beliebigen Beispiel zu beginnen:
\[
\vartheta_{00}(u+v)\,\vartheta_{11}(u-v) = A\,\vartheta_{01}(v)\,\vartheta_{10}(v) + B\,\vartheta_{00}(v)\,\vartheta_{11}(v),
\]
und wenn man $v=0$ setzt:
\[
A\,\vartheta_{01}\vartheta_{10} = \vartheta_{00}(u)\,\vartheta_{11}(u),
\]
daraus erhält man $B$ durch Vertauschung von $u$ mit $v$. So sind die drei folgenden Formeln abgeleitet:

\weq{5}{\begin{aligned}\vartheta_{01}\vartheta_{10}\,\vartheta_{00}(u+v)\,\vartheta_{11}(u-v)\\ {}= \vartheta_{00}(u)\vartheta_{11}(u)\vartheta_{01}(v)\vartheta_{10}(v) - \vartheta_{01}(u)\vartheta_{10}(u)\vartheta_{00}(v)\vartheta_{11}(v),\end{aligned}}
\weq{6}{\begin{aligned}\vartheta_{00}\vartheta_{01}\,\vartheta_{10}(u+v)\,\vartheta_{11}(u-v)\\ {}= \vartheta_{10}(u)\vartheta_{11}(u)\vartheta_{00}(v)\vartheta_{01}(v) - \vartheta_{00}(u)\vartheta_{01}(u)\vartheta_{10}(v)\vartheta_{11}(v),\end{aligned}}
\weq{7}{\begin{aligned}\vartheta_{00}\vartheta_{10}\,\vartheta_{01}(u+v)\,\vartheta_{11}(u-v)\\ {}= \vartheta_{01}(u)\vartheta_{11}(u)\vartheta_{00}(v)\vartheta_{10}(v) - \vartheta_{00}(u)\vartheta_{10}(u)\vartheta_{01}(v)\vartheta_{11}(v),\end{aligned}}
die sich unmittelbar verifizieren lassen, indem man erst $u$, dann $v=0$ setzt. Hieraus folgen die drei anderen durch Vermehrung von $u$ um
\[
\frac{1}{2}, \quad \frac{\omega}{2}, \quad \frac{1}{2}:
\]
% [sic] The source repeats 1/2 as the third shift.
\weq{8}{\begin{aligned}\vartheta_{01}\vartheta_{10}\,\vartheta_{01}(u+v)\,\vartheta_{10}(u-v)\\ {}= \vartheta_{01}(u)\vartheta_{10}(u)\vartheta_{01}(v)\vartheta_{10}(v) + \vartheta_{00}(u)\vartheta_{11}(u)\vartheta_{00}(v)\vartheta_{11}(v),\end{aligned}}
\weq{9}{\begin{aligned}\vartheta_{00}\vartheta_{01}\,\vartheta_{00}(u+v)\,\vartheta_{01}(u-v)\\ {}= \vartheta_{00}(u)\vartheta_{01}(u)\vartheta_{00}(v)\vartheta_{01}(v) - \vartheta_{10}(u)\vartheta_{11}(u)\vartheta_{10}(v)\vartheta_{11}(v),\end{aligned}}
\weq{10}{\begin{aligned}\vartheta_{00}\vartheta_{10}\,\vartheta_{00}(u+v)\,\vartheta_{10}(u-v)\\ {}= \vartheta_{00}(u)\vartheta_{10}(u)\vartheta_{00}(v)\vartheta_{10}(v) + \vartheta_{01}(u)\vartheta_{11}(u)\vartheta_{01}(v)\vartheta_{11}(v).\end{aligned}}

Es lassen sich diese Formeln, deren Gesamtheit mit dem Namen des Additionstheorems bezeichnet wird, in mannigfacher Weise verallgemeinern, wovon das folgende als Beispiel dienen möge.

Sind $u$, $v$ zwei Variable, so sind die vier Produkte
\[
\begin{array}{ll}
\vartheta_{00}(u)\,\vartheta_{00}(u+v), & \vartheta_{01}(u)\,\vartheta_{01}(u+v), \\
\vartheta_{10}(u)\,\vartheta_{10}(u+v), & \vartheta_{11}(u)\,\vartheta_{11}(u+v),
\end{array}
\]
%CONT>>
